\documentclass[12pt,a4paper]{article}
\usepackage[utf8]{vntex, inputenc}
\usepackage{amsmath, amssymb, anyfontsize, array, color,  enumerate, enumitem, fancyhdr, float, graphicx, hyperref, lastpage, listings, mathtools, mips, natbib, pifont, tikz}
\usepackage[open,openlevel=0]{bookmark}
\usepackage[framemethod=tikz]{mdframed}

\usetikzlibrary{calc}
\newcommand\HRule{\rule{12cm}{1pt}}

\pagestyle{fancy}
\fancyhead{} % clear all header fields
\fancyhead[L]{
	\color{blue}
	\begin{tabular}{rl}
		\begin{picture}(25,15)(0,0)
		\put(0,-8){\includegraphics[width=8mm, height=8mm]{BK.png}}
		\end{picture}
		\begin{tabular}{l}
			\textbf{\bf \ttfamily Ho Chi Minh University of Technology}\\
			\textbf{\bf \ttfamily Faculty of Computer Science \& Engineering}
		\end{tabular}
	\end{tabular}
}
\fancyhead[R]{
	\begin{tabular}{l}
		\tiny \bf \\
		\tiny \bf
	\end{tabular}
}
\fancyfoot{} % clear all footer fields
\fancyfoot[L]{\scriptsize \ttfamily Computer Architecture Week 4}
\fancyfoot[R]{\scriptsize \ttfamily Page {\thepage}/\pageref{LastPage}}
\renewcommand{\headrulewidth}{0.3pt}
\renewcommand{\footrulewidth}{0.3pt}

\definecolor{dkgreen}{rgb}{0,0.6,0}
\definecolor{gray}{rgb}{0.5,0.5,0.5}
\definecolor{mauve}{rgb}{0.58,0,0.82}

\lstset{frame=tb,
	language=[MIPS]Assembler,
	aboveskip=3mm,
	belowskip=3mm,
	showstringspaces=false,
	columns=flexible,
	basicstyle={\small\ttfamily},
	numbers=none,
	numberstyle=\tiny\color{gray},
	keywordstyle=\color{blue},
	commentstyle=\color{dkgreen},
	stringstyle=\color{mauve},
	breaklines=true,
	breakatwhitespace=true,
	tabsize=3,
	numbers=left,
	stepnumber=1,
	numbersep=1pt,
	firstnumber=1,
	numberfirstline=true
}

\begin{document}
\begin{titlepage}
  \begin{tikzpicture}[remember picture, overlay]
    \draw[line width = 2pt] ($(current page.north west) + (1in,-1in)$) rectangle ($(current page.south east) + (-1in,1in)$);
  \end{tikzpicture}

  \begin{center}

    % Upper part of the page
    \textbf{\fontsize{12pt}{1pt}\selectfont HO CHI MINH CITY UNIVERSITY OF TECHNOLOGY}\\[0.5cm]
    {\fontsize{13pt}{1pt}\selectfont Faculty of Computer Science \& Engineering}\\[0.5cm]
    \begin{figure}[H]
      \centering
      \includegraphics[width=1.7in,height=1.7in]{BK.png}
    \end{figure}

    % Title
    \HRule\\[0.5cm]
    { \textbf{\fontsize{25pt}{1pt}\selectfont Computer Architecture Week 4}}\\[0.4cm]

    \HRule\\[0.8cm]
    \begin{minipage}{0.545\textwidth}
      \begin{flushleft}
        \textbf{Authors:}\\
        \begin{tabular}{l l}
          Lưu Nguyễn Hoàng Minh & 1952845 \\
          Vũ Anh Nhi            & 1952380 \\
          Nguyễn Chính Khôi     & 1952793 \\
          Nguyễn Hoàng          & 1952255 \\
        \end{tabular}
      \end{flushleft}
    \end{minipage}
    \begin{minipage}{0.4\textwidth}
      \begin{flushright}
        \textbf{Class:}\\
        CO2008\_CC02\\
        \textbf{Lecturer:}\\
        Kiều Đỗ Nguyên Bình\\

      \end{flushright}
    \end{minipage}

    \vfill

    % Bottom of the page
    \vspace{2cm}
    {\large} %{\large \today}
  \end{center}
\end{titlepage}


\pdfbookmark[0]{Question 1}{ques1}
\section*{Question 1:}
Write a MIPS program with the following requirements:
\begin{enumerate}
  \item Request an integer number from users.
  \item If the number is positive, repeat step 1. Otherwise, print sum of all integer numbers that the program has read from users.
\end{enumerate}

\begin{mdframed}[hidealllines=true,backgroundcolor=magenta!10]
  \begin{lstlisting}

    .data:
    msg: .asciiz "Enter an integer: "
    msg1: .ascii "The entered number is not positive, the sum is: "

    # Start program code
    .text:

    .globl main

  main:
    # i : s0, sum : s1
    addi	$s0, $zero, 1 	# int i = 1
    addi	$s1, $zero, -1 	# int sum = -1

  WHILE:
    slt		$t0, $zero, $s0	# while (i > 0)
    bne		$t0, 1, WEXIT

    add		$s1, $s1, $s0 	# sum += i
    li		$v0, 4					# cout << "Enter an integer: "
    la		$a0, msg				#
    syscall								#
    li		$v0, 5				# cin >> i
    syscall								#
    move	$s0, $v0			#

    j WHILE

  WEXIT:
    # Print string msg1
    li		$v0, 4
    la		$a0, msg1
    syscall

    # Print sum
    li		$v0, 1
    move	$a0, $s1
    syscall

    # Exit
    li		$v0, 10
    syscall

  \end{lstlisting}
\end{mdframed}

\pdfbookmark[0]{Question 2}{ques2}
\section*{Question 2:}
Implement the following C code by using MIPS code. Assume that b and c are 10 and 5, respectively while input variable is read from keyboard. Print value of a to the terminal.
\begin{mdframed}[hidealllines=true,backgroundcolor=magenta!10]
  \begin{lstlisting}

  switch(input){
    case 0: a = b + c; break;
    case 1: a = b - c; break;
    case 2: a = c - b; break;
    deafult: printf("please input an another integer numbers\n"); break;
  }

  \end{lstlisting}
\end{mdframed}

\begin{mdframed}[hidealllines=true,backgroundcolor=magenta!10]
  \begin{lstlisting}

    .data:
    msg: .asciiz "Enter an integer: "
    print: .asciiz "a = "
    break: .asciiz "Please input other integer numbers"

    # Start program code
    .text:

    .globl main
  main:
    addi    $s1, $zero, 10         # b = 10
    addi    $s2, $zero, 5          # c = 5

    # Print string msg
    li      $v0, 4
    la      $a0, msg
    syscall

    # Get input and save
    li      $v0, 5
    syscall
    move    $s0, $v0

    add     $t0, $s1, $s2           # b + c
    sub     $t1, $s1, $s2           # b - c
    sub     $t2, $s2, $s1           # c - b

    beq     $s0, $t0, case0
    beq     $s0, $t1, case1
    beq     $s0, $t2, case2
    j default

  case0:
    # Print string print
    li      $v0, 4
    la      $a0, print
    syscall

    # Print a
    li      $v0, 1
    move    $a0, $s0
    syscall

    j Exit

  case1:
    # Print string print
    li      $v0, 4
    la      $a0, print
    syscall

    # Print a
    li      $v0, 1
    move    $a0, $s0
    syscall

    j Exit

  case2:
    # Print string print
    li      $v0, 4
    la      $a0, print
    syscall

    # Print a
    li      $v0, 1
    move    $a0, $s0
    syscall

    j Exit

  default:
    # Print string break
    li      $v0, 4
    la      $a0, break
    syscall

  Exit:
    # Exit
    li      $v0, 10
    syscall


  \end{lstlisting}
\end{mdframed}

\pdfbookmark[0]{Question 3}{ques3}
\section*{Question 3:}
Write a MIPS program with the following requirements.
\begin{enumerate}
  \item Declare an integer array with 10 synthetic data elements.
  \item Read an integer number from users.
  \item Find in the array if the integer read from user exists in the array or not. Print the position of the integer number in the array if found; otherwise tell users that the number does not exist in the array.
\end{enumerate}

\begin{mdframed}[hidealllines=true,backgroundcolor=magenta!10]
  \begin{lstlisting}

    .data:
    a: .word 1, 2, 3, 4, 5, 6, 7, 8, 9, 10
    msg: .asciiz "Enter an integer: "
    msg1: .asciiz "Position of the integer: "
    msg2: .asciiz "The integer does not exist in the array"

    # Start program code
    .text:

    .globl main

  main:
    li $s7, 40 	# array size
    li $s1, 0		# int i = 0
    li $s2, 0		# int pos = 0
    la $s3, a
    li $a0, 0

    # Print string msg
    li $v0, 4
    la $a0, msg
    syscall

    # Get input and save
    li $v0, 5
    syscall
    move $s0, $v0

  WHILE:
    slt $t0, $s1, $s7	# while (i < n)
    bne $t0, 1, WEXIT

    add $t1, $s1, $s3		# index [i]
    lw $t1, 0($t1)		# get a[i]
    beq $t1, $s0, TRUE
    addi $s1, $s1, 4
    addi $s2, $s2, 1
    j WHILE
  WEXIT:
    # Print string msg2
    li      $v0, 4
    la      $a0, msg2
    syscall
  j EXIT
  TRUE:
    # Print string msg1
    li      $v0, 4
    la      $a0, msg1
    syscall

    # Print pos
    li      $v0, 1
    move    $a0, $s2
    syscall
  EXIT:
    # Exit
    li      $v0, 10
    syscall

    \end{lstlisting}
\end{mdframed}


\pdfbookmark[0]{Question 4}{ques4}
\section*{Question 4:}
Given the following leaf procedure in ANSI C
\begin{mdframed}[hidealllines=true,backgroundcolor=magenta!10]
  \begin{lstlisting}

  void swap(int v[], int k)
  {
    int temp;
    temp = v[k];
    v[k] = v[k+1];
    v[k+1] = temp;
  }

  \end{lstlisting}
\end{mdframed}
Assume that the \(\$a0\) register will store the base address of the v array while the \(\$a1\) register keeps value k.
The array v consists of 10 elements in integer and is pre-defined in the data section.
\begin{enumerate}
  \item Write a main program the receive value k from user, check the value k and call the procedure swap if possible.
  \item Watch the \(\$ra\) register before and after the \(jal\) and \(jr\) instructions are executed.
\end{enumerate}
\begin{mdframed}[hidealllines=true,backgroundcolor=magenta!10]
  \begin{lstlisting}

    .data:
    a: .word 1, 2, 3, 4, 5, 6, 7, 8, 9, 10
    msg: .asciiz "Enter an integer: "
    msg1: .asciiz "New array is: "
    msg2: .asciiz " "
    msg3: .asciiz "Not possible"

    # Start program code
    .text:

    .globl main

  main:
    li $s7, 40 	# array size
    li $s1, 0	#int i = 0
    la $s2, a
    li $a0, 0

    # Print string msg
    li      $v0, 4
    la      $a0, msg
    syscall

    # Get input and save
    li      $v0, 5
    syscall
    move    $s0, $v0

    li $t1, 9
    slt $t0, $s0, $t1
    beq $t0, 0, NOTPOSSIBLE

    move $a0, $s2
    move $a1, $s0
    jal SWAP

    # Print string msg1
    li      $v0, 4
    la      $a0, msg1
    syscall
  WHILE:
    slt $t0, $s1, $s7	# while (i < n)
    bne $t0, 1, WEXIT

    add $t1, $s1, $s2	# index [i]
    lw $t1, 0($t1)		# get a[i]

    # Print a[i]
    li      $v0, 1
    move    $a0, $t1
    syscall
    # Print string msg2
    li      $v0, 4
    la      $a0, msg2
    syscall

    addi $s1, $s1, 4
    j WHILE
  WEXIT:
    j EXIT

  NOTPOSSIBLE:
    # Print string msg3
    li      $v0, 4
    la      $a0, msg3
    syscall
    EXIT:
    # Exit
    li      $v0, 10
    syscall

  SWAP:
    addi $sp, $sp, -4
    sw $s0, 0($sp)

    add $a1, $a1, $a1
    add $a1, $a1, $a1
    add $t0, $a0, $a1
    lw $s0, 0($t0)
    lw $t1, 4($t0)
    sw $t1, 0($t0)
    sw $s0, 4($t0)
    jr $ra

    lw $s0, 0($sp)
    addi $sp, $sp, 4

  \end{lstlisting}
\end{mdframed}


\pdfbookmark[0]{Question 5}{ques5}
\section*{Question 5:}
Given the following factorial MIPS program in a recursive form (as in the slide)
\begin{mdframed}[hidealllines=true,backgroundcolor=magenta!10]
  \begin{lstlisting}

  fact:
    addi    $sp, $sp, -8    # adjust stack for 2 items
    sw      $ra, 4($sp)     # save return address
    sw      $a0, 0($sp)     # save argument
    slti    $t0, $a0, 1     # test for n < 1
    beq     $t0, $zero, L1
    addi    $v0, $zero, 1   # if so, result is 1
    addi    $sp, $sp, 8     # pop 2 items from stack
    jr      $ra             # and return

  L1:
    addi    $a0, $a0, -1    # else decrement n
    jal     fact            # recursive call
    lw      $a0, 0($sp)     # restore original n
    lw      $ra, 4($sp)     # and return address
    addi    $sp, $sp, 8     # pop 2 items from stack
    mul     $v0, $a0, $v0   # multiply to get result
    jr      $ra             # and return

    \end{lstlisting}
\end{mdframed}
\begin{enumerate}
  \item Type the above procedure and write a main program that call the above procedure with different n, where n is in the \(\$a0\) register. Watch the results.
  \item When n is 2, run the program step by step and watch the execution of instructions as well as the \(\$ra\) register and values store/load to/from the stack.
\end{enumerate}
\begin{mdframed}[hidealllines=true,backgroundcolor=magenta!10]
  \begin{lstlisting}

    .data:
    msg: .asciiz "Enter an integer: "
    print: .asciiz "Output is "

    # Start program code
    .text:

    .globl main
  main:
    # Print string msg
    li      $v0, 4
    la      $a0, msg
    syscall

    # Get input and save
    li      $v0, 5
    syscall
    move    $s0, $v0

    # Caller
    move    $a0, $s0
    jal     fact
    move    $s0, $v0

    # Print string print
    li      $v0, 4
    la      $a0, print
    syscall

    # Print a
    li      $v0, 1
    move    $a0, $s0
    syscall

    # Exit
    li      $v0, 10
    syscall

  fact:
    addi    $sp, $sp, -8    # adjust stack for 2 items
    sw      $ra, 4($sp)     # save return address
    sw      $a0, 0($sp)     # save argument
    slti    $t0, $a0, 1     # test for n < 1
    beq     $t0, $zero, L1
    addi    $v0, $zero, 1   # if so, result is 1
    addi    $sp, $sp, 8     # pop 2 items from stack
    jr      $ra             # and return

  L1:
    addi    $a0, $a0, -1    # else decrement n
    jal     fact            # recursive call
    lw      $a0, 0($sp)     # restore original n
    lw      $ra, 4($sp)     # and return address
    addi    $sp, $sp, 8     # pop 2 items from stack
    mul     $v0, $a0, $v0   # multiply to get result
    jr      $ra             # and return

    \end{lstlisting}
\end{mdframed}

\end{document}